\chapter*{Introduction}
\addcontentsline{toc}{chapter}{Introduction}
\markboth{Introduction}{Introduction}

L'étude de la propagation des ondes en milieu désordonné se situe à l'intersection de la mécanique des milieux continus et de la dynamique des réseaux discrets. Bien que le transport ondulatoire soit classiquement modélisé par des équations macroscopiques basées sur une hypothèse d'homogénéisation, ce cadre devient caduc en présence d'hétérogénéités à petite échelle. La rupture de l'invariance par translation induite par le désordre structurel génère de multiples diffusions élastiques. Les interférences qui en résultent modifient les propriétés de transport et conduisent, dans le cas d'un système unidimensionnel, à un confinement spatial de l'énergie connu sous le nom de localisation d'Anderson.

Ce stage a pour objectif d'étudier ces phénomènes d'inhibition du transport et de confronter les approches théoriques aux simulations numériques. Le travail s'appuie sur le modèle de la chaîne harmonique unidimensionnelle, pour lequel des environnements de simulation ont été développés sous Python. Ces outils permettent d'implémenter des fluctuations stochastiques sur les paramètres du système, tels que les masses ou les raideurs, afin d'explorer la dynamique dans des régimes où les solutions analytiques sont limitées.

La caractérisation du régime de localisation linéaire repose principalement sur le formalisme des matrices de transfert. Cette méthode permet d'extraire l'exposant de Lyapunov et de quantifier la longueur de localisation, évaluant ainsi l'atténuation spatiale des ondes. Parallèlement à cette approche spectrale stationnaire, la dynamique temporelle du transport est étudiée en simulant la propagation de paquets d'ondes. L'observation de l'atténuation de ces impulsions permet d'établir le lien avec les mécanismes de transfert d'énergie. La combinaison de ces méthodes vise à décrire la transition entre un régime de transport propagatif et le confinement induit par la structure du milieu.

Le présent rapport s'articule en plusieurs étapes. Nous dresserons d'abord un état de l'art sur la localisation ondulatoire et son étude en régime unidimensionnel. Nous détaillerons ensuite les fondements cinématiques du modèle de la chaîne ordonnée et désordonnée, ainsi que les cadres théoriques et numériques employés. Enfin, nous analyserons les résultats de nos simulations en les confrontant aux prédictions analytiques et aux lois d'échelle de la localisation.